% GERADO por tools/generation/gerar_secoes_latex.py a partir de
% artifacts/paper/secoes/09-conclusao.md. Nao editar a mao: a proxima geracao
% desfaz. Corrija o Markdown e regere.
À \textbf{RQ1}, a análise responde que o MPO permite desvios quando deixa aberta a leitura de relações, processos, identidades e termos: a formalização publicada assumiu nove compromissos incompatíveis, classificados pela Representation Theory quando aplicável e rastreados até esses pontos de abertura. Isso demonstra permissão em uma formalização documentada, não frequência entre formalizadores. À \textbf{RQ2}, a OntoMPO revisada demonstra, no cenário avaliado, capacidade de representar fatos, consultá-los e comparar cobertura conceitual; a evidência oferece base para avaliação posterior, não prova eficácia organizacional, aderência de observatórios reais ou implantação operacional.

Para Sistemas de Informação, o trabalho deixa três contribuições. A primeira é uma \textbf{base semântica verificável para auditar aderência}: ela torna possível explicitar em que a aderência consiste, porque cada conceito tem definição e compromisso ontológico explícitos, e permite comparar cobertura conceitual entre iniciativas no cenário avaliado, sem inferir aderência de observatórios reais. A segunda é \textbf{rastreabilidade de proveniência para prestação de contas}: com o ciclo de extração, transformação e carga representado como acontecimento, com participantes e com tempo, torna-se possível percorrer o caminho do conteúdo divulgado até a origem do dado e dizer quem respondeu por cada etapa no cenário, oferecendo base para prestação de contas que a legislação de acesso à informação cobra de quem publica \cite{L12527}. A terceira é o \textbf{procedimento}, reusável fora deste caso: reconstruir a formalização como linha de base, confrontar cada compromisso formal com o trecho que o originou, exigir de todo achado o duplo teste, a violação ontológica e a abertura no texto, classificá-lo por uma tipologia estabelecida e publicar o antes e o depois como artefatos reexecutáveis.

Com isso, o trabalho responde, no II GranDSI-Br, ao desafio 3: \textit{Eco(Sistemas\textsuperscript{2}) de Informação} \cite{araujo2025grandsi}. A integração entre iniciativas depende de um acordo conceitual que não pode ser auditado enquanto o modelo de referência permanecer em prosa.

Quatro direções dão continuidade ao trabalho. A primeira é um \textbf{estudo comparativo entre observatórios reais}, que substitua o cenário sintético por iniciativas da literatura e meça a divergência entre implantações. A segunda é o \textbf{alinhamento com ontologias de gerenciamento de projetos}, que ligaria o projeto observado ao mesmo projeto visto por quem o executa. A terceira ataca o limite de um único formalizador: \textbf{formalizações independentes} do mesmo modelo, por equipes que não se comuniquem, mostrariam se os pontos de abertura identificados aqui produzem divergência de fato, e não apenas se a permitem. A quarta é a integração com \textbf{modelos de linguagem} como direção, e não como resultado: uma ontologia com definições e compromissos explícitos é insumo para extração assistida sobre documentação de observatórios, e o que se oferece a essa agenda é o artefato.
